% Final
\section{Návrh funkcí}

% Final
V této sekci se budu věnovat podrobnější specifikaci funkcí navrhované aplikace. Nejprve se zaměřím na popis jednotlivých aktérů, kteří budou s aplikací interagovat. Následně na základě funkčních požadavků sepíšu seznam případů užití, které budou podrobněji specifikovat, jak mají jednotlivé části aplikace fungovat z pohledu uživatele. Nakonec pro jednotlivé případy užití sepíšu podrobnou specifikaci jejich scénářů v jazyce Gherkin. Tyto scénáře budou sloužit jako zadání pro implementaci, ale také jako podklad pro spouštění automatizovaných testů pomocí nástroje Cucumber. Jelikož scénáře přesně popisují, jak se aplikace skutečně chová, bude je také možné v budoucnu využít jako podklad pro AI agenty, kteří budou moci například generovat nové funkce, testy nebo například uživatelskou dokumentaci aplikace.

% Final
\subsection{Popis aktérů}

% Final
V rámci specifikace funkcí je třeba popsat tzv. aktéry, kteří reprezentují role, které mohou s aplikací interagovat a vykonávat v ní určité akce \cite[21-22]{UseCases}. Tito aktéři budou využiti jak v popisu případů užití, tak ve specifikaci funkcí pomocí jazyka Gherkin. Na základě požadavků aplikace jsem identifikoval následující aktéry.

% Final
\subsubsection{Uživatel}

% Final
Uživatel je osoba, která využívá aplikaci k učení se programování. Mezi jeho hlavní aktivity patří například zapisování si kurzů, učení se programování, opakování látky, interakce s ostatními uživateli a interakce s prvky gamifikace v aplikaci.

% Final
\subsubsection{Nepřihlášený uživatel}

% Final
Nepřihlášený uživatel je osoba, která v aplikaci není přihlášena a nemá tak přístup k většině funkcí aplikace. Může se však přihlásit nebo se registrovat v aplikaci.

% Final
\subsubsection{Premium uživatel}

% Final
Premium uživatel je osoba, která má aktivní předplatné nebo jeho zkušební verzi. Premium uživatel může využívat všechny funkce aplikace jako běžný uživatel. Kromě těchto akcí může vykonávat i další akce, ke kterým nemá běžný uživatel přístup.

% Final
\subsubsection{Korektor}

% Final
Korektor je osoba zodpovědná za kontrolu kvality výukového obsahu, tedy například za kontrolu lekcí, cvičení a vysvětlení látky. Funkce aplikace může využívat ve stejném rozsahu jako premium uživatel. Kromě toho může využívat podrobnější nahlašování chyb při učení.

% Final
\subsection{Specifikace případů užití}

% Final
Jako další krok při návrhu funkcí jsem se rozhodl sepsat případy užití. Případy užití představují způsob, jak specifikovat funkční chování systému z pohledu jeho aktérů. Každý případ užití popisuje konkrétní cíl aktéra a interakce se systémem, které k dosažení tohoto cíle vedou. Díky tomuto přístupu je tak možné detailněji formulovat požadavky a identifikovat okrajové případy, které je třeba při návrhu a implementaci ošetřit. Případy užití tak mohou být užitečné pro vývojáře, protože jim poskytnou detailnější popis toho, jak má daná část systému fungovat. Dále mohou být využity testery například jako podklad pro vytváření automatizovaných testů. Díky tomu, že bývají případy užití specifikovány pomocí přirozeného jazyka, mohou být čitelné i pro osoby, které nejsou technicky zaměřené. Celkově tak tvoří přesnější specifikaci systému, na základě které lze následně postupovat při návrhu a implementaci samotné aplikace. \cite{UseCases}

% Final
V průběhu let se vyvinulo několik způsobů, jak k případům užití přistupovat a jak je specifikovat. Základy specifikace případů užití položil Ivar Jacobson v roce 1986. Alistair Cockburn \cite{Cockburn2001WritingEffectiveUseCases} následně popsal metodiku pro psaní podrobnějších případů užití. Přesto však prosazoval názor, že by forma specifikace případů užití měla odpovídat potřebám a rozsahu konkrétního projektu. Další techniku specifikace případů užití také popsal například Martin Fowler, který prosazoval názor, že neexistuje jeden standardní způsob, jak specifikovat případy užití, a že formát a rozsah případů užití by měly být zvoleny dle potřeb konkrétního projektu. \cite[22-27]{NURUIDesignSteps}

% Final
V této práci jsem se tak rozhodl vycházet z práce Alistaira Cockburna \cite{Cockburn2001WritingEffectiveUseCases} a specifikovat případy užití způsobem, který bude užitečný pro vývoj této konkrétní aplikace. Pro každý případ užití tak budu specifikovat jeho identifikátor, název, popis, hlavního aktéra, hlavní scénář a případné alternativní scénáře. 

% Final
Pro specifikaci scénářů jsem se rozhodl využít jazyk Gherkin, který přináší několik zásadních výhod, jimž se budu věnovat v dalších částech této kapitoly. Scénáře se v jazyce Gherkin specifikují přímo v repozitáři aplikace v souborech s příponou \texttt{.feature} \cite{CucumberDocs}. Proto jsem se rozhodl v tomto textu popsat jednotlivé scénáře pouze stručným popisem a detailně je pak specifikovat přímo v repozitáři aplikace. Díky tomuto přístupu tak budu moci využít jak výhod specifikace případů užití, tak i výhod jazyka Gherkin.

% Final
Přehled specifikovaných případů užití lze nalézt v tabulce \ref{table:use-cases-summary}. V příloze \ref{chapter:use-case-list} lze pak nalézt podrobnou specifikaci všech případů užití.

% Final
\captionsetup{justification=centering}
\begin{longtable}{@{} p{\linewidth} @{}}
    \caption{Přehled případů užití aplikace}
    \label{table:use-cases-summary}\\
    \toprule
    Název \\
    \midrule
    \endfirsthead
    \toprule
    Název \\
    \midrule
    \endhead
    \nameref{sec:uc-1} \\
    \nameref{sec:uc-2} \\
    \nameref{sec:uc-3} \\
    \nameref{sec:uc-4} \\
    \nameref{sec:uc-5} \\
    \nameref{sec:uc-6} \\
    \nameref{sec:uc-7} \\
    \nameref{sec:uc-8} \\
    \nameref{sec:uc-9} \\
    \nameref{sec:uc-10} \\
    \nameref{sec:uc-11} \\
    \nameref{sec:uc-12} \\
    \nameref{sec:uc-13} \\
    \nameref{sec:uc-14} \\
    \nameref{sec:uc-15} \\
    \nameref{sec:uc-16} \\
    \nameref{sec:uc-17} \\
    \nameref{sec:uc-18} \\
    \nameref{sec:uc-19} \\
    \nameref{sec:uc-20} \\
    \nameref{sec:uc-21} \\
    \nameref{sec:uc-22} \\
    \nameref{sec:uc-23} \\
    \nameref{sec:uc-24} \\
    \nameref{sec:uc-25} \\
    \nameref{sec:uc-26} \\
    \nameref{sec:uc-27} \\
    \nameref{sec:uc-28} \\
    \nameref{sec:uc-29} \\
    \nameref{sec:uc-30} \\
    \nameref{sec:uc-31} \\
    \nameref{sec:uc-32} \\
    \nameref{sec:uc-33} \\
    \nameref{sec:uc-34} \\
    \nameref{sec:uc-35} \\
    \nameref{sec:uc-36} \\
    \nameref{sec:uc-37} \\
    \nameref{sec:uc-38} \\
    \nameref{sec:uc-39} \\
    \nameref{sec:uc-40} \\
    \nameref{sec:uc-41} \\
    \nameref{sec:uc-42} \\
    \nameref{sec:uc-43} \\
    \nameref{sec:uc-44} \\
    \nameref{sec:uc-45} \\
    \bottomrule
\end{longtable}

% Final
\subsection{Specifikace funkcí pomocí jazyka Gherkin}\label{sec:gherkin-specification}

% Final
Po sepsání případů užití jsem dále nové funkce přesněji specifikoval pomocí jazyka Gherkin \cite{CucumberDocs}. Gherkin je jednoduchý strukturovaný jazyk, který umožňuje pomocí přirozeného jazyka popsat chování systému v podobě textových scénářů. Díky přesně dané struktuře je následně možné tyto scénáře strojově zpracovat pomocí nástroje Cucumber\footnote{Dostupné na: \href{https://cucumber.io/}{https://cucumber.io/}} a ověřit, že se aplikace chová v souladu se specifikací pomocí automatizovaných testů. Celý proces je postavený na přístupu \textit{Behavior-Driven Development} \cite{BDD}, tedy na metodice vývoje, v rámci které je možné přesně specifikovat očekávané chování systému a následně vytvořit automatizované testy, které ověří, že se aplikace skutečně podle této specifikace chová. Tento přístup vznikl jako reakce na rozšířenou metodiku \textit{Test-Driven Development} s cílem pomoci vývojářům a testerům nahlížet na testování více z pohledu samotného chování aplikace, přesněji pochopit a definovat požadované chování a minimalizovat tak nejednoznačnosti v zadání a implementaci jednotlivých funkcí \cite{BDDHistory}.

% Final
Tento přístup má několik zásadních výhod. Zaprvé, samotná specifikace slouží jako přesná dokumentace toho, jak se současný systém skutečně chová. Díky přirozenému jazyku může být užitečná nejen testerům a vývojářům, ale i lidem, kteří nejsou technicky zaměření. Zadruhé, díky automatizovaným testům, které jsou spouštěny přesně dle jednotlivých scénářů, je možné průběžně ověřovat, že se aplikace skutečně chová přesně v souladu se specifikací. Zatřetí, přesné specifikace funkcí mohou sloužit jako podklady pro asistenty umělé inteligence, kteří mohou pomáhat vývojářům a testerům s implementací kódu, tvorbou dokumentace a vytvářením testů. Používání těchto asistentů tak může být díky přesné specifikaci a automatizovaným testům výrazně efektivnější. \cite{CucumberDocs}

% Final
Z výše zmíněných důvodů jsem se rozhodl využít jazyk Gherkin a nástroj Cucumber pro specifikaci funkcí a automatizaci testů s cílem zajistit efektivnější vývoj a kvalitnější implementaci samotných funkcí aplikace.

% Final
\subsubsection{Postup při specifikaci}

% Final
Funkce a jejich scénáře jsem specifikoval na základě definovaných případů užití a v souladu s dokumentací nástroje Cucumber \cite{CucumberDocs}. Specifikace jsem prováděl přímo v repozitáři frontendu aplikace v souborech s příponou \texttt{.feature}. Tyto soubory jsem následně využil jako podklad pro implementaci a dále pro automatizaci testů, kterou popíšu v kapitole testování.

% Final
Soubory ve formátu Gherkin typicky obsahují definici funkce \texttt{Feature}, jeden nebo více scénářů \texttt{Scenario} a jejich jednotlivé kroky. Kroky jsou definovány pomocí klíčových slov, kterými jsou například slova \texttt{Given} definující předpoklady, \texttt{When} definující interakci s aplikací, \texttt{Then} definující očekávaný výsledek a \texttt{And} umožňující definovat více kroků stejného typu za sebou. Příklad definice funkce pomocí jazyka Gherkin lze nalézt v ukázce kódu \ref{fig:feature-file-example}.

% Final
\begin{listing}[ht]
    \begin{minted}{gherkin}
@courses
Feature: Course and career path browsing

  Scenario: User opens a course from the courses page
    Given I am signed in
    And I am on the "courses-page" page
    Then I should see the "all-courses-section" section
    And I should see the "recommended-courses-section" section
    When I click on the "course-card" card
    Then I should be on the "course-page" page
    \end{minted}
    \caption{Ukázka specifikace funkce prohlížení kurzů pomocí jazyka Gherkin}
    \label{fig:feature-file-example}
\end{listing}
